\chapter{Introdução}

Este livro contém uma colecção de ensinamentos de Ajahn Sumedho dados a pessoas
que estão familiarizadas com as convenções do Budismo Theravāda e que têm alguma
experiência em meditação. A maioria dos capítulos são editados de palestras
proferidas durante retiros para leigos ou para os discípulos monásticos
(ordenados) de Ajahn Sumedho, portanto, requerem alguma atenção cuidadosa e são
melhor lidos em sequência.

Nos retiros monásticos de dois meses, Ajahn Sumedho desenvolve um tema a partir
dos ensinamentos do Buddha, ligando-o a outros aspectos do Dhamma, embelezando-o
com relatos das suas experiências pessoais, demonstrando a sua relevância para a
sociedade em geral, ou usando-o como uma exortação ao Sangha para viver de
acordo com a sua aspiração à iluminação. Embora não seja possível traduzir a
profundidade do tom e a variedade dessas palestras numa obra impressa,
acrescentada de breves exortações, directrizes e reflexões contemplativas mais
longas envolvidas com os cânticos que os monges e monjas recitam diariamente há
anos, pode sugerir a atmosfera e a perspectiva em que os ensinamentos são
oferecidos.

% REVISÃO: "Geração Dependente" — os capítulos dedicados ao tema (19–23) usam "Génese Dependente". Uniformizar o termo (recomenda-se "Génese Dependente") em todo o livro.
Em muitas dessas palestras, Ajahn Sumedho expõe a expressão singular budista de
``não-eu'' (\emph{anatta}). Ele afirma que esta é a forma do Buddha apontar para
a experiência da Realidade Última, que é o objectivo de muitas religiões.
Durante os retiros monásticos, Ajahn Sumedho frequentemente ensina a Geração
Dependente (\emph{paṭiccasamuppāda}) com base na abordagem de \emph{anatta}. A
Geração Dependente traça o processo pelo qual o sofrimento (\emph{dukkha}) é
composto pela ignorância (\emph{avijjā}) e é, inversamente, eliminado (ou
melhor, não gerado) pela cessação da ignorância. Assim como \emph{anatta},
não-eu, é a expressão da Verdade Suprema, Ajahn Sumedho sugere que a raiz da
ignorância é a ilusão de ``Eu''. É a sua tentativa de não aniquilar ou rejeitar
qualidades pessoais, mas sim indicar como o sofrimento (\emph{dukkha}) surge ao
tentar sustentar uma identidade denotada por corpo e mente.

% REVISÃO: aspas inconsistentes — usar ``Génese Dependente'' (aspas duplas, como no resto do texto) em vez de `Geração Dependente'.
% REVISÃO: vírgula a separar o sujeito do verbo ("A profundidade da ... Dependente, descreve-nos") — remover a vírgula após "Dependente"; frase de difícil leitura.
Esta identidade equivocada é o que a pessoa comum chama de ``mim''. Pode ser
detectada num estado latente como auto-consciência, como humor habitual da
mente, como vaidade, auto-criticismo, ou pode manifestar-se como actividade
egoísta corporal ou verbal. A profundidade da `Geração Dependente',
descreve-nos, como mesmo na sua forma mais passiva, essa visão errada cria
impulsos habituais (kamma) e atitudes, através das quais até mesmo um meditador
silencioso e bem-intencionado experimenta sofrimento. O Kamma vai do plano
psicológico ``interno'' até à realidade ``externa'' da acção. Este processo
habitual manifesta-se então em termos de corpo, fala ou mente; sendo todas essas
manifestações denominadas \emph{saṅkhāra}. Mesmo a acção moral baseada na
``visão egotista'', pode conduzir à ansiedade, dúvida, ``tristeza, lamento, dor,
angústia e desespero''. Esse é o significado do primeiro ``elo'' da Geração
Dependente ``\emph{avijjā paccayā saṅkhārā}'' ou ``as formações kammicas
dependem da ignorância''.

Na sua formulação mais completa, a Geração Dependente é expressa como:

\begin{quote}
  \raggedright
  Avijjā-paccayā saṅkhārā,
  saṅkhāra-paccayā viññāṇaṁ,
  viññāṇa-paccayā nāma-rūpaṁ,
  nāma-rūpa-paccayā saḷ-āyatanaṁ,
  saḷ-āyatana-paccayā phasso,
  phassa-paccayā vedanā,
  vedanā-paccayā taṇhā,
  taṇhā-paccayā upādānaṁ,
  upādāna-paccayā bhavo,
  bhava-paccayā jāti,
  jāti-paccayā jarā-maraṇaṁ soka-parideva-dukkha-domanass'upāyāsā sambhavanti.
  Evam-etassa kevalassa dukkhakkhandhassa samudayo hoti.
\end{quote}

% REVISÃO: "Isto trata-se de" é incorrecto (o sujeito de "tratar-se" é indeterminado). Corrigir para "Trata-se do surgimento de \emph{dukkha}." ou "Isto é o surgimento de \emph{dukkha}."
Isto trata-se do surgimento de \emph{dukkha}.

A cessação de \emph{dukkha} é então mapeada:

\begin{quote}
  \raggedright
  Avijjāya tveva asesa-virāga-nirodhā saṅkhāra-nirodho,
  saṅkhāra-nirodhā viññāṇa-nirodho,
  viññāṇa-nirodhā nāma-rūpa-nirodho,
  nāma-rūpa-nirodhā saḷ-āyatana-nirodho,
  saḷ-āyatana-nirodhā phassa-nirodho,
  phassa-nirodhā vedanā-nirodho,
  vedanā-nirodhā taṇhā-nirodho,
  taṇhā-nirodhā upādāna-nirodho,
  upādāna-nirodhā bhava-nirodho,
  bhava-nirodhā jāti-nirodho,
  jāti-nirodhā jarā-maraṇaṁ soka-parideva-dukkha-domanass'upāyāsā nirujjhanti.
  Evam-etassa kevalassa dukkhakkhandhassa nirodho hoti.

  \quoteRef{\href{https://suttacentral.net/sn12.35/en/bodhi}{SN 12.35, Avijjāpaccaya Sutta}}
\end{quote}

% REVISÃO: "o desejo depende da sensação; o desejo depende da sensação, a posse depende do desejo," — a vírgula após "depende do desejo" deveria ser ponto e vírgula, para manter a pontuação uniforme da enumeração.
% REVISÃO: incoerência entre os dois parágrafos — "parideva" é traduzido por "lamento" no primeiro e por "lamentação" no segundo (e "doença/morte" vs. lista distinta). Uniformizar a tradução dos termos da fórmula.
Em inglês, isto pode ser traduzido como:

\begin{quote}
  As formações habituais (kamma) dependem da ignorância; a consciência depende
  das formações habituais (kamma); o nome-forma (mentalidade-corporalidade)
  depende da consciência; as seis bases dos sentidos dependem do nome-forma; o
  contacto depende das seis bases dos sentidos; a sensação depende do contacto;
  o desejo depende da sensação; a posse depende do desejo, o devir depende da
  posse; o nascimento depende do devir; do nascimento dependem: a velhice, a
  doença e a morte, a tristeza, lamento, dor, angústia e desespero.

  Através de toda a cessação desta ignorância, cessam as formações habituais;
  através da cessação das formações habituais, cessa a consciência; através da
  cessação da consciência, cessa o nome-forma; através da cessação do
  nome-forma, cessam as seis bases dos sentidos; através da cessação das seis
  bases dos sentidos, cessa o contacto; através da cessação do contacto, cessa a
  sensação; através da cessação da sensação, cessa o desejo; através da cessação
  do desejo, cessa o apego; através da cessação do apego, cessa o devir; através
  da cessação do devir, cessa o nascimento; através da cessação do nascimento,
  cessam a velhice, a doença e a morte, a tristeza, a lamentação, a dor, a
  angústia e o desespero. Assim é a cessação de toda esta massa de sofrimento.
\end{quote}

Existem muitas formas de dependência em questão nesta análise. É útil lembrar
que \emph{paccayā} ``dependente de'' ou ``condições'' não significa
necessariamente ``gerar''. Por exemplo, poderia dizer-se ``andar depende das
pernas'', ou ``o gelo depende da água'', ou ``apanhar o comboio depende de
chegar à estação à hora certa'', ou mesmo ``a vista depende do não aparecimento
de objectos intervenientes''. Compreendendo isto, o contemplativo começa a
perceber que, assim como ``dependência emergente'' não significa necessariamente
``geração'', a ``cessação'' tão valorizada pelo Buddha não significa
necessariamente ``aniquilação''. Nesta vida, onde o Nibbāna deve ser
% REVISÃO: "formações-karmicas" — usar "kámmicas"/"kammicas" (de "kamma"), como em "formações kammicas" mais acima; evitar o hífen e a grafia "karmicas".
% REVISÃO: vírgula a separar o sujeito do verbo ("a identificação ... mentais, pode cessar") — remover a vírgula antes de "pode".
realizado, a mentalidade-corporalidade pode ``cessar'' -- ou seja, a
identificação com as formações-karmicas físicas e mentais, pode cessar de
modo a que a vida não seja mais vivida a partir do princípio do prazer/dor
ditado pelos sentidos. (\emph{nāma-rūpa-saḷāyatana-phassa-vedanā-taṇhā}). Com
este espírito, pode-se interpretar a sequência de uma forma mais fluida, por
exemplo:

% REVISÃO: "objeto" deveria ser "objecto" (grafia pré-Acordo, como no resto do livro).
Na medida em que (\emph{paccayā}) a mente não compreendeu (\emph{avijjā}) a
Verdade, impulsos habituais (\emph{saṅkhārā}) manifestam e condicionam
(\emph{paccayā}) a consciência para um modo discriminativo (\emph{viññāṇa}) que
opera em termos de (\emph{paccayā}) sujeito e objeto (\emph{nāma-rūpa}) mantidos
(\emph{paccayā}) para existir num lado ou outro das seis portas dos sentidos
(\emph{saḷāyatana}).

% FIXME: domanassa missing

Estas portas dos sentidos abrem-se dependendo (\emph{paccayā}) do contacto
(\emph{passā}) que pode suscitar (\emph{paccayā}) vários graus de sensação
(\emph{vedanā}). A sensação estimula (\emph{paccayā}) o desejo (\emph{taṇhā}) e,
de acordo com (\emph{paccayā}) o poder do desejo, a atenção mantém-se
(\emph{upādāna}) e assim se desenvolvem objectivos e obsessões pessoais
(\emph{bhava}) que dão (\emph{paccayā}) o ciclo de amadurecimento e
desaparecimento (\emph{jarā-maraṇa}) com o resultante sentimento de tristeza
(\emph{soka}) variando entre angústia (\emph{parideva}), depressão
(\emph{dukkha}) e esgotamento nervoso (\emph{upāyāsa}).

Quando a mente olha para o sentimento de perda e compreende a Verdade
\emph{(avijjā-nirodha)}, os impulsos habituais cessam \emph{(saṅkhāra-nirodha)}
e a consciência não mais é limitada pela discriminação dos mesmos
\emph{(viññāṇa-nirodha)}, de modo que a separação entre sujeito e objecto não
mais é sustentada \emph{(nāma-rūpa nirodha)} e não nos sentimos encurralados
dentro ou puxados para fora das seis portas dos sentidos
\emph{(saḷāyatana-nirodha)}. As portas dos sentidos abrem-se à reflexão, ao
invés de ficarem dependentes do contacto \emph{(phassa-nirodha)} e o impacto não
se imprime adentro na mente (\emph{vedanā-nirodha}). Portanto, existe liberdade
em relação ao desejo \emph{(taṇhā-nirodha)} e a atenção não fica presa
\emph{(upādāna-nirodha)} nem cresce para motivações egoístas
\emph{(bhava-nirodha)} que reforçam e se centram no ego. Quando nenhuma imagem
pessoal é criada, jamais poderá inflar nem ser destruída
\emph{(jarā-maraṇa-nirodha)}. Assim, nada há a perder, permanecendo uma sensação
de felicidade, elevação, alegria e serenidade
\emph{(soka-parideva-dukkha-domanassa-upāyāsā-nirodha)}.

Com a cessação de tal ponto de referência limitado à morte, passa-se a viver a
Verdadeira vida, a vida Santa, da qual Ajahn Sumedho tão evocativamente fala.

Embora muitas destas palestras tenham sido dadas para monásticos, a beleza do
Dhamma é estar disponível para aqueles que desejam ouvir. É com isto em mente
que este livro é oferecido gratuitamente.

\bigskip

{\raggedright
  Que todos os seres realizem a Verdade,\\
  Ven. Sucitto Bhikkhu\\
  Amaravati
\par}
